% XeLaTeX 编译两次。正文最多4页，参考文献和提示词附录另起页。
\documentclass{article}
\usepackage[preprint,nonatbib]{neurips_2026}
\usepackage[UTF8,fontset=fandol]{ctex}
\usepackage{graphicx,booktabs,tabularx,amsmath,xcolor,fancyvrb,hyperref}
\hypersetup{colorlinks=true,urlcolor=blue,linkcolor=black}
\makeatletter
\renewcommand{\@noticestring}{自然语言处理与大语言模型\quad 实验课1报告}
\makeatother
\linespread{1.15}
\title{论文引用数比较中的提示词探索：\\显式判据与位置去偏能否帮助4B思考模型？}
\author{姓名\quad 学号\\复旦大学}
\begin{document}
\maketitle
\begin{abstract}
本实验让Qwen3-4B（thinking）根据两篇论文的标题、摘要和日期判断哪篇引用数更高。我们依次实现了直接指令、检索式few-shot、结构化分步比较、多次采样投票加A/B交换去偏四种方案，在100条公开集上的准确率分别为68\%、69\%、72\%、73\%，最终方案达到公开集最高档门槛。误差分析表明：零样本存在明显的“偏向A”的位置偏置，交换位置投票能将其消除；few-shot改变了29题的答案却只净增1题；所有提升都在一个标准误（约4.5个百分点）以内；另有10题四种方法全部答错，主要源于模型对“领域热度”和“贡献类型”的误判以及提示词中关于发表时间的先验与数据不符。
\end{abstract}

\section{问题分析}
\paragraph{任务特点。}每道题给出同一时期、通常同一领域的两篇论文，要求判断记录时引用数更高的一篇。引用数取决于领域规模、贡献类型（方法、数据集、工具、综述还是个案研究）、成果的通用性以及积累引用的时间，这些信息只能从摘要中间接推断，没有确定答案，因此即使是人类专家也难以做到很高的准确率。我们在100条示例池上做了简单统计（未使用公开集标签）：答案A、B各占50\%；两篇论文日期差的中位数约3个月、最大不超过11个月，“较早发表的一篇胜出”的比例仅为49\%；“摘要较长的一篇胜出”为58\%，是一个弱信号。

\paragraph{研究问题与预期。}本实验关注的问题是：\emph{对一个4B规模的思考模型，“给出明确的比较维度”与“抵消位置偏置”能否带来稳定的提升？}我们的预期是：(1)小模型零样本时容易受选项顺序影响，交换A/B并投票可以抵消这种偏置；(2)明确列出引用数的驱动因素能让模型的推理更聚焦；(3)few-shot示例的标签信息有限（只给出答案，不给理由），作用可能不大。

\paragraph{整体方案。}所有方法共用同一套论文格式（标题、日期、摘要，空白字符规整为单个空格）和答案提取逻辑：要求模型最后一行输出\texttt{Answer: A/B}，用正则表达式取最后一次出现的答案；若最终回答中提取不到，再从思考文本的末尾600个字符中提取；仍失败则重试，所有调用都失败时返回兜底答案A。单次请求的总长度控制在19000字节以内（框架上限为20000字节）。各步骤通过常量\texttt{STEP}切换，提示词集中写在\texttt{prompts.py}中，并自动同步到\texttt{solution.py}。

\section{提示词调优过程}
\paragraph{方法1：直接指令（zero-shot）。}system提示把模型设定为“预测论文引用数的科学计量学专家”，user提示给出两篇论文，只要求输出一行答案。这是基线，每题调用1次。

\paragraph{方法2：检索式few-shot。}从100条示例池中按词重叠（去停用词后的Jaccard相似度）检索与当前题最相似的4条示例，其中答案为A、B的各2条，以避免标签偏置；示例按相似度升序排列，使最相似的示例离问题最近。示例以多轮对话形式给出（user轮是论文对，assistant轮是\texttt{Answer: X}），超过字节上限时先丢弃最不相似的示例。

\paragraph{方法3：结构化分步比较。}在方法2的示例检索基础上，更换system与user提示：system中列出引用数的四个主要驱动因素（领域规模、贡献类型与可复用性、结论的强度与新颖性、积累时间，并说明几个月的日期差远不如领域和贡献类型重要）；user中要求模型按“领域与社区规模→贡献类型→结论的重要性与通用性→日期差的影响”四步比较后再做决定。日期差由代码预先算好，以“Paper A比Paper B早发表N天”的形式写入提示，避免小模型自己做日期运算出错（上下文工程）。

\paragraph{方法4：多次采样投票 + A/B交换去偏（最终方案）。}以方法3为单次判断，每题调用4次，其中第2、4次交换两篇论文的位置，并把答案映射回原顺序；按多数票决定，2:2平票时追加一次原顺序调用，仍平票则取第一个有效答案。实际共调用417次，说明有17题出现了平票。

\begin{table}[h]
\centering
\caption{不同方法在同一份100条公开集上的运行结果。“预测A数”为100题中输出A的数量（真实标签中A、B各50个）。}
\label{tab:main}
\small
\begin{tabularx}{\linewidth}{lXrrrrr}
\toprule
方法 & 主要思路 & 正确数/100 & 调用数 & 无效返回 & 预测A数 & 总Token\\
\midrule
方法1 & 直接指令，单次调用 & 68 & 100 & 0 & 58 & 14.5万\\
方法2 & 相似示例4-shot（A、B各2条） & 69 & 100 & 0 & 49 & 37.8万\\
方法3 & 判据+分步比较+日期差+4-shot & 72 & 100 & 0 & 52 & 44.5万\\
方法4 & 方法3×4次投票，交换A/B & \textbf{73} & 417 & 0 & 47 & 187.4万\\
\bottomrule
\end{tabularx}
\end{table}

\paragraph{结果与分析。}表\ref{tab:main}中四种方法都没有无效返回，也没有因长度截断而失败，说明答案提取逻辑足够稳健。主要观察如下：

\begin{itemize}
\item \textbf{零样本存在位置偏置。}方法1预测A 58次，对真实答案为A的题召回38/50，对真实答案为B的题只召回30/50。加入A、B平衡的示例后，预测A的次数降到49；方法4通过交换位置投票进一步降到47，对A、B两类的召回分别为35和38，偏置基本消除。
\item \textbf{few-shot的作用接近随机扰动。}方法2与方法1有29题答案不同，其中方法2对15题、方法1对14题，净增只有1题。示例只提供“答案”而不提供“理由”，4B模型很难从4个例子中归纳出可迁移的判断规律；同时prompt长度增加到原来的约5倍。
\item \textbf{显式判据带来较大的单步提升}（69→72）。在模型的思考文本中可以看到它确实按“领域—贡献类型—重要性—日期”的顺序逐项比较，判断依据比方法1更具体。
\item \textbf{投票的边际收益很小。}方法4与方法3只有15题答案不同（8:7，净增1题），调用数和Token却是方法3的约4倍。17题出现2:2平票，说明这些题上模型的判断本身就不稳定，多数票只能在两个近似随机的答案间做选择。
\end{itemize}

\paragraph{统计显著性。}在100题上，70\%左右准确率的标准误约为$\sqrt{0.7\times0.3/100}\approx4.6$个百分点。方法1与方法4的答案有29题不同，其中方法4对17题、方法1对12题，McNemar检验$\chi^2=(|17-12|-1)^2/29\approx0.55$，$p\approx0.46$，差异并不显著。因此上述排序只能视为趋势，私有集上的结果可能会有±5个百分点的波动。

\paragraph{最终选择。}我们最终提交方法4，理由是：(1)它在公开集上准确率最高（73\%）；(2)它对A、B两类答案的召回最均衡，不依赖选项顺序，在分布未知的私有集上更稳健；(3)417/500的调用量留有余量，不会因预算耗尽导致漏答。代价是运行时间约11分钟，是方法1的7倍。

\clearpage
\section{实验心得}
\begin{figure}[h]
\centering
\IfFileExists{figures/run.png}{\includegraphics[width=\linewidth,height=5cm,keepaspectratio]{figures/run.png}}{\fbox{\parbox[c][3cm][c]{0.92\linewidth}{\centering 放入自己的运行与评测截图\\展示结果排行榜、完成条数和调用预算\\保存为figures/run.png}}}
\caption{真实运行记录（\texttt{runs/reports/}中的结果排行榜）。四次完整运行分别为68、69、72、73题正确，与表\ref{tab:main}一致。}
\end{figure}

\paragraph{方法之间的互补性。}四种方法两两之间的答案一致率在71\%--85\%之间，方法3与方法2、方法3与方法4都是85\%。至少被一种方法答对的题有90道，只有10道题（P012、P017、P039、P044、P052、P062、P074、P081、P083、P091）四种方法全部答错。但对四种方法的答案简单取多数只得到71题正确，低于方法4，说明各方法的“对”并不集中在可以通过投票挑出来的地方。

\paragraph{代表性错误案例。}
\begin{itemize}
\item \textbf{P039}：A是深度学习训练基准\emph{TBD: Benchmarking and Analyzing Deep Neural Network Training}（2018年3月），B是高斯过程软件库\emph{GPyTorch}（2018年9月），真实答案为B。模型抓住了system提示中“benchmark被引更多”这一条，同时认为A“早半年发表、有时间优势”，于是选了A；但它忽略了GPyTorch作为通用工具库的长期被引潜力。
\item \textbf{P044}：A是用现成生成模型发现视觉模型缺陷（2022年8月），B是\emph{Score Jacobian Chaining}，把2D扩散模型提升到3D生成（2022年12月），真实答案为B。模型低估了“扩散模型+3D生成”在当时的热度，又一次把较早的日期当作优势。
\item \textbf{P012 / P081}：两题都是天体物理领域内部的比较（矮星系化学组成 vs.\ 黑洞双星；伽马射线望远镜触发策略 vs.\ 尘埃磁取向）。同一领域内的细分方向热度差异，需要相当具体的领域知识，4B模型往往只能给出含糊的“两者都是小众方向”的判断，最后的选择接近随机。
\end{itemize}

\paragraph{哪些改动没有达到预期。}(1)\textbf{日期先验与数据不符。}方法3的提示中写了“更早的论文有优势”，并专门提供日期差，但示例池中较早一篇胜出的比例仅为49\%；P039、P044都是模型因为“早发表”而选错。这条先验来自常识，没有先用数据验证，是本实验最明显的设计失误。(2)\textbf{few-shot示例没有理由。}只给答案的示例几乎没有带来增益，反而让prompt长度增加约5倍。(3)\textbf{同一提示重复投票多样性不足。}方法3与方法4的答案高度相关，投票主要抵消了位置偏置，而没有纠正系统性的判断错误。

\paragraph{总结与改进方向。}本实验最有价值的发现有三点：小模型零样本存在可测量的位置偏置，交换A/B投票能有效消除；明确列出判断维度比堆叠示例更有效；在100题规模上，几个百分点的差异很容易被噪声淹没，结论必须结合显著性来看。局限在于所有调参都只在100条公开集上进行，存在过拟合风险，而私有集的满分门槛（75\%）更高。后续可尝试：(1)删除或弱化日期先验，改为“日期差基本不影响结果”；(2)先让模型独立给每篇论文的引用潜力打分，再比较分数，以减少成对比较中的顺序干扰；(3)固定5次奇数票，并混合不同风格的提示词，提高投票的多样性；(4)给few-shot示例配上简短的判断理由；(5)把示例池也当作验证集，并在选例时排除当前题本身，使用200题比较方案以降低噪声。

\clearpage
\section*{参考文献}
\begin{enumerate}
\small
\item T.~Brown et al. Language Models are Few-Shot Learners. \emph{NeurIPS}, 2020.
\item J.~Wei et al. Chain-of-Thought Prompting Elicits Reasoning in Large Language Models. \emph{NeurIPS}, 2022.
\item X.~Wang et al. Self-Consistency Improves Chain of Thought Reasoning in Language Models. \emph{ICLR}, 2023.
\item J.~Liu et al. What Makes Good In-Context Examples for GPT-3? \emph{DeeLIO Workshop}, 2022.
\item L.~Zheng et al. Judging LLM-as-a-Judge with MT-Bench and Chatbot Arena. \emph{NeurIPS Datasets and Benchmarks}, 2023.
\item Qwen Team. Qwen3 Technical Report. \emph{arXiv:2505.09388}, 2025.
\end{enumerate}

\appendix
\section{最终选定的提示词}
最终方案为方法4（\texttt{STEP = 4}），它每次调用的单次判断使用方法3的提示词。以下内容与提交的\texttt{solution.py}中\texttt{PROMPTS}区块一致，为便于排版对长行做了手动换行，花括号内为动态字段。

\paragraph{论文格式（所有方法共用）。}
\begin{Verbatim}[fontsize=\small,frame=single]
PAPER_TEMPLATE:
Paper {label}
Title: {title}
Date: {date}
Abstract: {abstract}

PAIR_TEMPLATE:
{paper_a}

{paper_b}
\end{Verbatim}

\paragraph{system提示（COT\_SYSTEM）。}
\begin{Verbatim}[fontsize=\small,frame=single]
You are an expert in scientometrics. Your job is to predict which of two
papers has more citations. Citation counts are driven mainly by: (1) the
size and activity of the research community of the topic; (2) how broadly
useful the contribution is (new general methods, datasets, benchmarks,
software tools, surveys and foundational theory are cited far more than
narrow case studies or incremental variants); (3) the strength and novelty
of the claims; (4) the time a paper has had to accumulate citations (an
earlier paper has an advantage, but a gap of a few months matters much less
than topic and contribution type).
\end{Verbatim}

\paragraph{user提示（COT\_USER）及日期说明。}
\begin{Verbatim}[fontsize=\small,frame=single]
Predict which paper had received MORE citations at the time the data was
recorded.

{pair}

Extra context: {date_note}

Compare the papers step by step:
1. Field and community size of each paper (is it a hot, large area or a
   small niche?).
2. Type of contribution (method / dataset / benchmark / tool / survey /
   theory / application / case study) and how widely reusable it is.
3. Significance and generality of the results as claimed in the abstract.
4. Effect of the publication date gap.
Then weigh these factors and decide.

The LAST line of your reply must be exactly "Answer: A" or "Answer: B".

date_note（由代码计算）：
  Paper {older} was published {gap} days earlier than Paper {newer}.
  或（同一天）：Both papers were published on the same day.

few-shot示例的assistant轮（FEWSHOT_ANSWER）：
Answer: {answer}
\end{Verbatim}

\paragraph{消息构造与答案汇总。}
\begin{enumerate}
\item \textbf{选例}：对当前题和示例池中的每条示例，取两篇论文的标题与摘要，小写化、去掉停用词和长度不超过2的词，计算词集合的Jaccard相似度；分别取答案为A和B的最相似的2条，共4条，按相似度升序排列。
\item \textbf{消息}：\texttt{[system: COT\_SYSTEM]}，然后对每条示例依次加入\texttt{user: COT\_USER(示例)}和\texttt{assistant: Answer: X}，最后是\texttt{user: COT\_USER(当前题)}。总长度超过19000字节时，从最不相似的示例开始逐条删除。
\item \textbf{投票}：每题调用4次，第2、4次交换\texttt{paper\_a}与\texttt{paper\_b}（示例保持不变，日期说明随之重算），并把答案翻转回原顺序；平票时追加1次原顺序调用，仍平票取第一个有效答案；全部失败时返回A。
\item \textbf{提取}：在最终回答中依次尝试\texttt{Answer: X}、单独一行的\texttt{X}、加粗的\texttt{**X**}三种模式，取最后一次匹配；若都没有匹配，再在思考文本的末尾600个字符中提取。
\end{enumerate}
\end{document}
